% Appendix: the solutions to the practice exam, from exam/practice_exam_solutions.md, with the
% points of each part and what earns part of them: the source's italic notes are \credit, its
% "Fälla:" notes \trap.
%
% One deliberate difference: question 6 c) asks for a complete program, and the book prints no
% program as a solution. Its marking criteria and its trap are kept; the program is not.

% The solutions' headings are labelled pans:q:N.
\renewcommand{\answerprefix}{pans}

\chapter{Lösningsförslag till övningsprovet}
\label{app:practice:answers}

Lösningarna följer övningsprovet i bilaga~\ref{app:practice}. Vid varje deluppgift står poängen och
vad som ger delpoäng. Bedömningsprinciperna, med metod och följdfel, står i bilaga~\ref{app:exam}.
Programmen i fråga 6 a) och b) har körts i simulatorn.

% ==========================================================================================
\answersection{1}{Talsystem och koder}{8}

\answerpart{a}{2}
Ettorna har vikterna 128, 64, 16, 4 och 2: \code{128 + 64 + 16 + 4 + 2 = 214}. Grupperna \code{1101}
och \code{0110} är \code{D} och \code{6}: \textbf{214} och \textbf{\code{0xD6}}.
\credit{1 p per rätt omvandling med uträkning.}

\answerpart{b}{2}
173: 128 ryms, 45 kvar; 32 ryms, 13 kvar; 8 ryms, 5 kvar; 4 ryms, 1 kvar; 1 ryms. Ettor på vikterna
128, 32, 8, 4 och 1: \textbf{\code{1010 1101}}. Grupperna \code{1010} och \code{1101} är \code{A}
och \code{D}: \textbf{\code{0xAD}}.
\credit{1 p per rätt omvandling. Division med 2 och avläsning av resterna är lika bra.}

\answerpart{c}{2}
\begin{textdrawing}
  minnessiffror: 1 11
                   1011 0100      0xB4 = 180
                 + 0110 1010    + 0x6A = 106
                 -----------
                 1 0001 1110      286
\end{textdrawing}

\noindent Bit~5 ger en minnessiffra in i bit~6, bit~6 en in i bit~7, och bit~7 en som går ut ur
registret.

Registret får de åtta låga bitarna, \textbf{\code{0x1E}}, och \textbf{en minnessiffra går ut} ur
registret, eftersom \code{180 + 106 = 286} inte ryms i åtta bitar: \code{286 - 256 = 30 = 0x1E}.
\credit{1 p för en korrekt uppställd addition, 1 p för \code{0x1E} och minnessiffran.}

\answerpart{d}{2}
BCD skriver varje siffra för sig med fyra bitar: \textbf{\code{0101 1001}}. ASCII-koderna för
tecknen \code{'5'} och \code{'9'} är \textbf{\code{0x35}} och \textbf{\code{0x39}}.
\credit{1 p för BCD, 1 p för ASCII-koderna.}
\trap{59 binärt, \code{0011 1011}, är inte BCD och ger 0~p på BCD-delen.}

% ==========================================================================================
\answersection{2}{Grindar och boolesk algebra}{8}

\answerpart{a}{2}
\begin{narrowtable}{@{}cccccc@{}}
  A & B & C & \code{A'B} & \code{BC'} & X \\
  \midrule
  0 & 0 & 0 & 0 & 0 & 0 \\
  0 & 0 & 1 & 0 & 0 & 0 \\
  0 & 1 & 0 & 1 & 1 & 1 \\
  0 & 1 & 1 & 1 & 0 & 1 \\
  1 & 0 & 0 & 0 & 0 & 0 \\
  1 & 0 & 1 & 0 & 0 & 0 \\
  1 & 1 & 0 & 0 & 1 & 1 \\
  1 & 1 & 1 & 0 & 0 & 0 \\
\end{narrowtable}
\credit{2 p för en helt rätt tabell, 1 p för högst ett fel.}

\answerpart{b}{3}
\begin{textdrawing}
X = AB + AB' + A'B
  = A(B + B') + A'B        bryt ut A
  = A · 1 + A'B            komplement: B + B' = 1
  = A + A'B                identitet
  = A + B                  förenklingslagen: A + A'B = A + B
\end{textdrawing}

\noindent Det förenklade uttrycket är en \textbf{OR-grind}.
\credit{1 p för att slå ihop de två första termerna till A, 1 p för förenklingslagen, 1 p för OR.}

\noindent Ett svar som stannar vid \code{A + A'B} får 1~p: det saknar förenklingslagen, och pekar
inte ut någon enda grind.

\answerpart{c}{3}
\begin{enumerate}
  \item \textbf{\code{A'B'}}, med De Morgan. \worth{1}
  \item Med De Morgan är \code{A + B = (A'B')'}. Inversen av en variabel är en NAND-grind med båda
    ingångarna kopplade till variabeln: \code{(AA)' = A'}. Nätet blir tre NAND-grindar: en som
    bildar \code{A'}, en som bildar \code{B'}, och en tredje med \code{A'} och \code{B'} in, som
    bildar \code{(A'B')' = A + B}. \worth{2}
    \credit{1 p för ett korrekt nät, 1 p för att visa varför det stämmer.}
\end{enumerate}

% ==========================================================================================
\answersection{3}{Karnaughdiagram}{6}

\answerpart{a}{2}
\begin{narrowtable}{@{}ccccc@{}}
  \thead{AB \textbackslash{} CD} & \thead{00} & \thead{01} & \thead{11} & \thead{10} \\
  \midrule
  \textbf{00} & 1 & & & 1 \\
  \textbf{01} & & 1 & 1 & \\
  \textbf{11} & & 1 & 1 & \\
  \textbf{10} & 1 & & & 1 \\
\end{narrowtable}
\credit{1 p för rätt Graykod på båda axlarna, 1 p för rätt placerade ettor.}
\trap{axlarna i binär ordning, \code{00, 01, 10, 11}, ger 0~p på axlarna, men följdfelet bedöms i
b).}

\answerpart{b}{3}
Två grupper om fyra:
\begin{itemize}
  \item \textbf{De fyra hörnen}, mintermerna 0, 2, 8 och 10. Hörnen är grannar eftersom diagrammet
    går runt både uppifrån och ned och från vänster till höger. Där är \code{B} = 0 och \code{D} =
    0: termen \textbf{\code{B'D'}}.
  \item \textbf{De fyra mittrutorna}, mintermerna 5, 7, 13 och 15. Där är \code{B} = 1 och \code{D}
    = 1: termen \textbf{\code{BD}}.
\end{itemize}

\begin{console}
X = B'D' + BD
\end{console}
\credit{1 p för varje rätt grupp med rätt term, 1 p för det samlade uttrycket.}

\noindent Den som missar att hörnen är en grupp och skriver hörnen som två par,
\code{A'B'D' + AB'D' + BD}, får 1~p.

\answerpart{c}{1}
\code{X} beror bara på \textbf{B och D}, och är 1 när de är lika: en \textbf{XNOR}-grind. A och C
påverkar inte \code{X} alls.

% ==========================================================================================
\answersection{4}{Analys av ett kontaktnät}{6}

\answerpart{a}{2}
Två parallellkopplingar i serie. Den första är \code{S1 + S2}, den andra \code{S1 + S3'}, eftersom
S3 är brytande:

\begin{console}
H1 = (S1 + S2)(S1 + S3')
\end{console}
\credit{1 p per rätt parallellkoppling, inklusive \code{S3'}.}

\answerpart{b}{2}
\begin{narrowtable}{@{}cccccc@{}}
  S1 & S2 & S3 & \code{S1 + S2} & \code{S1 + S3'} & H1 \\
  \midrule
  0 & 0 & 0 & 0 & 1 & 0 \\
  0 & 0 & 1 & 0 & 0 & 0 \\
  0 & 1 & 0 & 1 & 1 & 1 \\
  0 & 1 & 1 & 1 & 0 & 0 \\
  1 & 0 & 0 & 1 & 1 & 1 \\
  1 & 0 & 1 & 1 & 1 & 1 \\
  1 & 1 & 0 & 1 & 1 & 1 \\
  1 & 1 & 1 & 1 & 1 & 1 \\
\end{narrowtable}
\credit{2 p för en helt rätt tabell, 1 p för högst ett fel. Följdfel från a) bedöms här.}

\answerpart{c}{2}
Med den distributiva lagen \code{A + BC = (A + B)(A + C)}, läst baklänges:

\begin{console}
H1 = (S1 + S2)(S1 + S3') = S1 + S2S3'
\end{console}

\noindent Nätet: den slutande kontakten S1 \textbf{parallellt med} en seriekoppling av den slutande
kontakten S2 och den brytande kontakten S3, och sedan lampan. Tre kontakter i stället för fyra, och
bara ett kontaktblock på S1.
\credit{1 p för det förenklade uttrycket, 1 p för ett rätt ritat nät.}

\noindent Kontrollera med tabellen: \code{S1 + S2S3'} är 1 på raderna 010, 100, 101, 110 och 111,
samma som i b).

% ==========================================================================================
\answersection{5}{Mikrodatorns uppbyggnad}{4}

\answerpart{a}{2}
CPU:n med ALU (räknar och jämför), register (håller värdena som räknas med), styrenhet (avkodar
instruktionerna och styr resten) och programräknare (pekar på nästa instruktion). Programminnet
håller programmet, dataminnet variablerna och stacken, och in- och utportarna förbinder
mikrodatorn med omvärlden. Bussarna, adressbussen och databussen, förbinder CPU:n med minnena och
portarna.
\credit{1 p för ett blockschema med alla delar på plats, 1 p för att beskriva vad de gör.}

\answerpart{b}{1}
Programräknaren innehåller \textbf{adressen till nästa instruktion} som ska hämtas. Vid
\code{rcall} sparas adressen till instruktionen efter \code{rcall} på stacken, och
programräknaren får subrutinens adress, så att nästa instruktion hämtas därifrån.

\answerpart{c}{1}
Programmet ligger i \textbf{flashminnet}, som behåller sitt innehåll utan ström. Stacken ligger i
\textbf{SRAM}, som töms när strömmen bryts.

% ==========================================================================================
\answersection{6}{Assemblerprogrammering}{8}

\answerpart{a}{3}
Uppmätt i simulatorn: \textbf{\code{r16} = \code{0x68}}, \textbf{\code{r17} = \code{0x04}},
\textbf{\code{r18} = \code{0x2C}}, och \textbf{\code{C} = 1} efter \code{add}.
\begin{itemize}
  \item \code{and r17, r16}: \code{0000 1111 AND 1011 0100 = 0000 0100} = \code{0x04}.
  \item \code{lsl r16}: \code{1011 0100} skiftas ett steg åt vänster och blir \code{0110 1000} =
    \code{0x68}. Bit~7 hamnar i \code{C}, men \code{C} skrivs över senare av \code{add}.
  \item \code{add r18, r19}: \code{200 + 100 = 300}, som inte ryms i åtta bitar. Registret får
    \code{300 - 256 = 44} = \code{0x2C}, och minnessiffran hamnar i \code{C}.
\end{itemize}
\credit{1 p för \code{r16} och \code{r17}, 1 p för \code{r18}, 1 p för \code{C} och förklaringen:
åtta bitar räcker till 255.}

\answerpart{b}{2}
\begin{enumerate}
  \item Loopen körs \textbf{4 gånger}, en gång för varje värde på \code{r17}: 4, 3, 2 och 1.
    \code{PORTB} får värdena 1, 2, 4 och 8, och är \textbf{\code{0x08}} när loopen är klar.
    \code{r16} är då \code{0x10}, men det skrivs aldrig till porten.
  \item Ett varv är \code{out} 1, \code{lsl} 1, \code{dec} 1 och \code{brne} 2 cykler, alltså 5;
    det sista varvet 4. Loopen tar \code{4 · 5 - 1 = 19} cykler, och \code{19 · 62,5 ns ≈ 1,19 µs}.
    Simulatorn mäter 19 cykler.
\end{enumerate}
\credit{1 p per deluppgift.}
\trap{\code{PORTB} = \code{0x10}, genom att räkna med \code{lsl} efter den sista \code{out}, ger
0~p på den delen.}

\answerpart{c}{3}
\noanswer{Programmet står inte i boken. Så här bedöms det:}
\credit{1 p för initieringen (utgång, ingång och pull-up), 1 p för att läsa rätt bit och ta hänsyn
till att knappen är aktivt låg, 1 p för en loop som tänder och släcker rätt. Ett program som läser
hela porten med \code{in}, inverterar med \code{com} och maskar med \code{andi} är lika bra.}
\trap{ett program som tänder lysdioden när biten är 1 lyser när knappen är släppt, och förlorar
poängen för aktivt låg.}
